\chapter{أوجه اختلاف الأحاديث (الجمع والترجيح)}

\section{اختلاف التنوع (صيغ التشهد والقراءات)}
يتابع الإمام الشافعي رحمه الله مناقشة الأحاديث التي ظاهرها الاختلاف، ويبين كيف يتعامل معها الفقيه. ومن ذلك اختلاف الصحابة في ألفاظ "التشهد". فقد روى ابن مسعود، وابن عباس، وعمر بن الخطاب، وعائشة، وجابر، وابن عمر رضي الله عنهم صيغاً للتشهد تختلف في بعض الألفاظ بالزيادة والنقصان والتقديم والتأخير.

والحقيقة أن هذا من "اختلاف التنوع" وليس اختلاف تضاد. فالنبي صلى الله عليه وسلم نوّع في العبادة تيسيراً وتوسعة على الأمة، فكل هذه الصيغ صحيحة ومجزئة، وكلها تعظيم لله عز وجل. وقد اختار الشافعي تشهد ابن عباس رضي الله عنهما (التحيات المباركات الصلوات الطيبات لله...) لأنه أوسع وأكثر لفظاً، مع إقراره بصحة الباقي وعدم تعنيفه لمن أخذ بغيره مما ثبت عن النبي صلى الله عليه وسلم. 

\subsection{القرآن أُنزل على سبعة أحرف}
ومما يشبه اختلاف التنوع، اختلاف القراءات القرآنية. فقد استدل الشافعي بحديث عمر بن الخطاب حين سمع هشام بن حكيم يقرأ سورة الفرقان على غير ما أقرأه النبي صلى الله عليه وسلم، فلما احتكما إليه قال: «إن هذا القرآن أُنزل على سبعة أحرف فاقرؤوا ما تيسر منه». فهذا تنويع وتوسعة ورأفة من الله بالأمة لاختلاف لهجاتهم، ما لم يكن في الاختلاف إحالة وتغيير للمعنى. وإذا جاز هذا الاختلاف في لفظ القرآن للتوسعة، فجوازه في السنن والأذكار (كالتشهد) من باب أولى.

\section{الجمع والترجيح في أحاديث الربا}
ذكر الشافعي اختلاف الرواية على وجه آخر يحتاج إلى جمع أو ترجيح. فالأحاديث المتواترة (كحديث أبي سعيد الخدري، وأبي هريرة، وابن عمر، وعبادة بن الصامت) تحرم "ربا الفضل"، وتنص على أنه لا يجوز بيع الذهب بالذهب، والفضة بالفضة إلا مثلاً بمثل ويداً بيد دون زيادة.

ولكن جاء في المقابل حديث أسامة بن زيد رضي الله عنهما أن النبي صلى الله عليه وسلم قال: «إنما الربا في النسيئة» (أي ربا الأجل). وظاهره حصر الربا في النسيئة فقط وإباحة ربا الفضل.

\textbf{طرق توجيه حديث أسامة:}
\begin{enumerate}
    \item \textbf{مسلك الجمع:} يُحمل حديث أسامة على اختلاف الأصناف (كبيع الذهب بالفضة، أو التمر بالقمح)، فهنا يجوز التفاضل، ويكون الربا المحرم هو النسيئة (التأخير) فقط. أو أن أسامة رضي الله عنه أدرك الجواب ولم يسمع السؤال الذي قيد هذه الإجابة بمسألة معينة.
    \item \textbf{مسلك الترجيح:} إن تعذر الجمع، صرنا إلى الترجيح؛ فالذين رووا تحريم ربا الفضل جماعة كبار أحفظ وأقدم صحبة (كعثمان وعبادة وأبي هريرة وابن عمر)، وحديث الجماعة أولى بالقبول من حديث الواحد.
\end{enumerate}

\section{اختلاف الأحاديث في وقت صلاة الفجر}
ومن الاختلاف الظاهري ما ورد في وقت صلاة الفجر؛ ففي حديث عائشة وزيد بن ثابت وسهل بن سعد: أن النبي صلى الله عليه وسلم كان يصلي الصبح بغلس (أي في الظلمة). وفي المقابل حديث رافع بن خديج: «أسفروا بالفجر فإنه أعظم للأجر». 

\textbf{توجيه الشافعي للتعارض:}
بيّن الشافعي أننا إذا احتجنا للترجيح، نرجح حديث التغليس والمبادرة للصلاة لأسباب:
\begin{itemize}
    \item أنه أشبه وأقرب لظاهر القرآن في المبادرة للخيرات ﴿حَافِظُوا عَلَى الصَّلَوَاتِ وَالصَّلَاةِ الْوُسْطَىٰ﴾، فمن قدمها في أول وقتها كان أولى بالمحافظة.
    \item أنه مروي عن عدد أكبر وأحفظ من الصحابة.
    \item أن الخلفاء الراشدين (أبا بكر وعمر وعثمان وعلي) كانوا يدخلون فيها بغلس.
\end{itemize}
أما حديث الإسفار، فيُحمل على التأكد من طلوع الفجر الصادق (المعترض في الأفق) وعدم التعجل قبل دخول الوقت، أو أن الإسفار يكون في الانصراف منها بسبب إطالة القراءة فيها.

\section{اختلاف الحال في النهي عن استقبال القبلة للغائط والبول}
ومن الاختلاف الذي يرجع لتغير الحال: حديث أبي أيوب الأنصاري في النهي المطلق: «لا تستقبلوا القبلة ولا تستدبرها لغائط أو بول». يقابله حديث ابن عمر: «رأيت رسول الله صلى الله عليه وسلم على لبنتين مستقبلاً بيت المقدس لحاجته» (أي مستدبراً الكعبة).

\textbf{طريقة الجمع بينهما:} 
حمل الشافعي حديث النهي على قضاء الحاجة في "الصحراء" والفضاء المفتوح؛ لسعتها ولئلا تُقذر القبلة أو تُكشف العورات تجاه المصلين. وحمل رخصة الاستدبار (أو الاستقبال) على "البنيان والمنازل"؛ لضيقها ووجود الستر فيها. فابن عمر حدث بما رأى في البنيان، وأبو أيوب حدث بما سمع من النهي العام، والفقيه يجمع بينهما بالتفريق بين حال الصحراء وحال البنيان.

\section{خاتمة في أدب الاختلاف}
إن الخلاف إما أن يكون اختلاف أفهام وتنوع، كصيغ التشهد وقراءة القرآن، فهذا يُعذر فيه المخالف ولا يُعنف. وإما أن يكون اختلاف تضاد ومصادمة للنصوص، كمن يحلل ما حرم الله، أو من يعارض الدليل الواضح بالهوى والآراء المحدثة؛ فهذا لا يُقبل، والكثرة لا تبرر الباطل ﴿وَإِن تُطِعْ أَكْثَرَ مَن فِي الْأَرْضِ يُضِلُّوكَ عَن سَبِيلِ اللَّهِ﴾. 
وقد ضرب الشافعي أروع الأمثلة في أدب الخلاف، حين رد على شيخه سفيان بن عيينة في مسألة نسخ أحاديث، فرد عليه بالأدلة والتاريخ دون تنقيص، ليعلمنا أن العلم مبني على الدليل الشرعي وليس على التقليد الأعمى.